\section{Background}
\label{sec:background}

The challenge of governing recursive spawning in multi-agent AI systems---where one agent spawns child agents that may themselves spawn further descendants---sits at the intersection of three distinct research traditions: agent lifecycle management, distributed-systems provenance, and hierarchical task decomposition. This section reviews each in turn, then positions the BX3 Framework as the unifying substrate that addresses the critical gap among them: immutable parent-chain traceability across recursive agent spawning.

% ==============================================================================
\subsection{Agent Lifecycle Management in Multi-Agent Systems}
\label{sec:lifecycle}

Modern AI agent platforms do not generally manage the full lifecycle of their constituent agents with the rigor applied to traditional software processes. An agent may be invoked, may spawn children, may merge results, and may terminate---but these transitions are typically implicit, undocumented, and invisible to downstream auditors. This gap has driven substantial recent work.

\bayeretali{bayer2025identity} observe that the shift from deterministic software to autonomous agents acting on external services introduces a fundamental identity-management challenge: traditional enterprise identity infrastructure (SSO, SCIM) was not designed for actors whose behavior is non-deterministic and whose action surfaces span multiple domains. They argue that standardized interoperable identity profiles are necessary to maintain auditable agent identities from creation through retirement, and they highlight MCP (Model Context Protocol) as a leading framework for connecting language models to external tools in multi-agent workflows. Their contribution emphasizes that per-agent accountability records must be established at provisioning time and maintained across delegation events.

\auditableai{auditability2025} found that auditability must be treated as a first-class system requirement for multi-agent AI trustworthiness. Their five dimensions of agent auditability---action recoverability, lifecycle coverage, policy checkability, responsibility attribution, and evidence integrity---provide a concrete evaluation framework for any spawning infrastructure. Critically, they demonstrate that no agent system can be truly accountable without auditability, and that evidence integrity (tamper-evident logging of delegation events) is a prerequisite for the other four dimensions.

\provcagent{provcagent2025} propose unified provenance tracking for agentic workflows, demonstrating that provenance graphs must capture not just data transformations but agent spawn events, delegation boundaries, and intent propagation across the lifecycle. Their \texttt{prov:wasDerivedFrom} and \texttt{prov:actedOnBehalfOf} extensions to the W3C PROV model provide a formal vocabulary for modeling parent-child spawning relationships. \trusttrack{trusttrack2025} add a protocol-stack approach to verifiable agents, showing that trust-native agent systems require provenance tracking as a first-class runtime property alongside authentication and authorization.

A separate thread examines lifecycle governance from the perspective of authorization scope. The Human Delegation Provenance (HDP) Protocol \citep{draft-helixar-hdp} introduces a cryptographic token scheme that binds a human authorization event to a session identifier and records each delegation hop as a signed, append-only entry. HDP enables offline verification of the full delegation chain using only the issuer's public key, without central registries or trusted third parties. While originally proposed for human-to-orchestrator delegation, its append-only chain model maps directly to machine-to-machine parent-child spawning in hierarchical agent systems.

\aip{aip2025} introduce Invocation-Bound Capability Tokens (IBCTs), which fuse identity, authorization attenuation, and provenance into a single token chain. Their evaluation shows that verification overhead remains sub-millisecond even for chains of depth five, with token size scaling linearly at approximately 340--380 bytes per delegation block. This performance envelope is critical: any provenance mechanism whose verification cost grows superlinearly with chain depth becomes impractical for real-time agentic systems with deep spawning trees.

\mooreetal{moore2025hmas} provide a five-axis taxonomy for hierarchical multi-agent systems covering decomposition topology, delegation authority models, communication patterns, failure handling, and lifecycle governance---offering a useful conceptual map for comparing recursive spawning architectures along structural dimensions.

% ==============================================================================
\subsection{Accountability and Provenance in Distributed Systems}
\label{sec:provenance}

Provenance in distributed systems is not a new problem, but AI agents amplify it in two specific ways: the rate at which spawning and delegation events occur can exceed human review capacity, and the semantic content of delegation (goals, constraints, context) is richer than in traditional service calls. These conditions make traditional audit logging insufficient and have driven new protocol design.

\sentinelagent{sentinelagent2025} formalize delegation integrity with the Delegation Chain Calculus (DCC), which defines seven properties of sound delegation chains---six deterministic (authority narrowing, policy preservation, forensic reconstructibility, cascade containment, scope-action conformance, output schema conformance) and one probabilistic (intent preservation). SentinelAgent implements these properties at runtime via a non-LLM Delegation Authority Service, achieving a 100\% true-positive rate with 0\% false-positive on DelegationBench under black-box conditions. Their three-point verification lifecycle (pre-execution intent checking, at-execution scope enforcement, post-execution output validation) provides a concrete engineering pattern for enforcing delegation integrity as a first-class runtime property. These seven properties serve as functional requirements against which any recursive spawning substrate must be evaluated.

\loka{loka2025} introduce a decentralized framework for trustworthy multi-agent ecosystems organized around a Universal Agent Identity Layer (UAIL), intent-centric communication, and a Decentralized Ethical Consensus Protocol. While oriented toward ethical alignment, their UAIL provides the identity substrate required for accountable spawning: agents carry verifiable identities across interactions, enabling post-hoc reconstruction of which agent performed which action under whose authority. This property is essential for the forensic audit trails required in regulated industries such as healthcare, finance, and critical infrastructure.

\chainofconsciousness{chainofconsciousness2026} present a cryptographic protocol for verifiable AI agent provenance using hash chains, where each agent's spawn event and outputs are sealed into a Merkle chain. Their approach guarantees that any tampering with the delegation record is detectable, trading dynamic adaptability for cryptographic non-repudiation of the spawn lineage.

\aaudittrack{aaudittrack2025} extend this with a decentralized interaction-graph framework for AI security and accountability, modeling agent spawning as nodes in a graph and applying graph-theoretic accountability metrics to determine responsible parties across delegation chains. Their framework demonstrates that structural properties of the spawn graph (depth, branching factor, cross-branch communication) directly affect the system's capacity to assign post-hoc accountability.

\layeredmutability{layeredmutability2026} address the tension between persistence and governance in self-modifying agents, showing that layered mutability models---where some properties of an agent are immutable and others are mutable under governance---provide a principled way to balance adaptability with auditability. Their work informs the design of spawning systems where child agents inherit immutable parent pointers while retaining mutable execution scopes.

Collectively, this body of work establishes that accountability and provenance in spawning-based systems require structural guarantees, not merely policy enforcement. The spawning architecture itself must make accountability computationally tractable, or the audit burden grows faster than the system's productive capacity.

% ==============================================================================
\subsection{Fixed vs. Mutable Delegation Chains}
\label{sec:delegation-chains}

A critical architectural question for any spawning-based system is whether the delegation chain is fixed at creation or mutable during execution. The distinction has profound implications for both security and auditability.

A \emph{fixed delegation chain} establishes the full parent-child lineage at spawn time; children know their ancestors, and any subsequent re-delegation requires explicit new chain construction. This model is exemplified by the immutable parent pointers in the BX3 Framework's recursive spawning architecture (Section \ref{sec:framework}), where each spawned Worksheet contains a hard-coded pointer to the parent's Purpose Layer. The property is structural, not policed: the chain cannot be broken by context loss or network failure because the parent reference is architecturally indestructible. Every descendant node can trace its accountability lineage to a human anchor regardless of spawning depth.

A \emph{mutable delegation chain} permits dynamic re-delegation, where an agent may redirect work to a different child based on load, capability, or intermediate results. The Recursive Delegation Engine (RDE) framework \citep{intelligent2025} operationalizes this with a directed graph $G = (A, E)$ where agents may update their delegation edges based on learned utility signals, using multi-armed bandit strategies (epsilon-greedy, UCB, Thompson Sampling) to optimize routing. TDAG \citep{tdag2024} similarly generates specialized subagents dynamically during execution, arguing that static agent assignments introduce brittleness in complex, multi-step tasks. Their dynamic ToA (Tree of Agents) permits runtime re-branching as task decomposition evolves.

\agenthive{agenthive2025} propose dynamic, decentralized delegation as a first-class primitive for LLM-powered multi-agent systems, showing that delegation graphs must adapt to intermediate results and capability signals. They demonstrate that decentralized delegation improves throughput in horizontally scalable workloads but at the cost of increased provenance complexity.

The tradeoff is between adaptability and auditability. Mutable chains can route around failures and exploit capability diversity, but they complicate forensic reconstruction: if an agent can re-delegate to any peer at runtime, the \emph{actual} delegation path may differ from the \emph{intended} one. Fixed chains are more restrictive but guarantee that the audit trail reflects the design-time intent.

The HDP Protocol \citep{draft-helixar-hdp} and IBCT architecture \citep{aip2025} both implement a hybrid variant: the delegation chain itself is append-only and immutable, but agents may create new child tokens that extend the chain. This preserves the audit trail's completeness while allowing dynamic subtask creation. Each new hop is cryptographically signed, making the extended chain fully verifiable without requiring real-time communication.

The BX3 Framework resolves this tension architecturally: the parent chain is immutable (guaranteeing traceability), but the Bounds Engine still permits arbitrary task decomposition and subtask execution within the inherited scope (preserving adaptability within accountability constraints). Critically, BX3's immutability is a \emph{structural property of the spawn event}, not a runtime policy that can be overwritten by context loss or an overwhelmed model. This is the key distinction from policy-based and cryptographically-hybrid approaches.

% ==============================================================================
\subsection{Hierarchical Task Decomposition in AI}
\label{sec:task-decomp}

Hierarchical task decomposition is the mechanism by which a complex goal is segmented into a chain or tree of simpler subgoals, each assigned to a specialized agent or processing unit. It is the structural foundation of any recursive spawning-based multi-agent architecture.

\intelligent{intelligent2025} propose an adaptive delegation framework that explicitly models task decomposition as a parent-child authority transfer. Each delegation event includes not only the sub-goal but the scope of authority, accountability boundaries, and intent---making the decomposition itself part of the audit record. Their framework includes mechanisms for trust-building and provenance propagation between delegator and delegatee, enabling reliable collaboration in agentic networks. Crucially, they demonstrate that dynamic reuse and recovery from failures are achievable within a hierarchical structure when authority transfers are treated as first-class objects rather than side effects.

\firmhive{firmhive2025} demonstrate hierarchical task decomposition at scale in the FIRMHIVE framework for autonomous firmware analysis. FIRMHIVE builds a runtime Tree of Agents (ToA) where delegation is treated as an executable primitive: agents autonomously decompose goals into per-agent executable subgoals rather than receiving static task assignments. Their approach achieves approximately $16\times$ more reasoning steps and $2.3\times$ broader file inspection compared to flat agent baselines, demonstrating that hierarchical decomposition is not merely a governance mechanism but an effective scalability mechanism.

\lcm{lcm2025} introduce Lossless Context Management (LCM), a deterministic memory architecture for LLM-based agents that replaces model-written loops with engine-managed, operator-level recursion. LCM's dual-state memory (immutable store plus active context) preserves parent-chain lineage across decomposition steps, enabling drill-down into any prior decision without recomputation. Their \texttt{llm\_map} and \texttt{agentic\_map} primitives implement bounded, terminating task decomposition as declarative engine operations rather than heuristic model behavior. In benchmarks across context lengths from 32K to 1M tokens, LCM-augmented agents outperform alternatives that rely on model-driven recursion loops.

\reactree{reactree2025} present ReAcTree, a hierarchical framework for long-horizon task planning using LLM-based agents arranged in a dynamic tree. Control flow nodes coordinate strategy; agent nodes reason, act, and expand the tree. Each agent node uses goal-specific episodic memory and shares environment observations through working memory, enabling context-aware decomposition that adapts to intermediate results. ReAcTree demonstrates that hierarchical decomposition plus memory-augmented agents improves robustness and efficiency in embodied long-horizon tasks, achieving approximately $61\%$ goal success versus $31\%$ for flat ReAct baselines on WAH-NL.

\agentorchestra{agentorchestra2025} present AgentOrchestra, a hierarchical multi-agent framework that manages agent specialization and coordination through structured spawning, showing that hierarchical topologies outperform flat ones on general-purpose task solving by enabling role specialization while maintaining central coordination.

Collectively, this body of work establishes that hierarchical task decomposition is both necessary for scalable multi-agent systems and compatible with strong auditability guarantees when the decomposition itself is governed as a first-class artifact rather than an emergent side effect.

% ==============================================================================
\subsection{The BX3 Framework as Substrate for Immutable Parent Chains}
\label{sec:bx3-substrate}

The BX3 Framework, introduced in full in Section \ref{sec:framework}, provides a specific architectural answer to the fragmentation described above. Its three functional layers---Purpose Layer (intent, judgment, accountability), Bounds Engine (bounded reasoning, constraint interpretation), and Fact Layer (deterministic enforcement, forensic auditability)---are defined by the properties they must maintain, not by the nature of the actor occupying them.

Applied to recursive spawning, the BX3 Framework produces a specific structural guarantee: when a parent system spawns a child, the child receives a Worksheet containing the complete system state and a hard-coded pointer to the parent's Purpose Layer. This pointer cannot be overwritten by context loss or network conditions. It is a structural property of the spawn event, not a runtime policy. The result is an immutable parent chain: every descendant node can trace its accountability lineage to a human anchor, regardless of the depth or complexity of the spawning tree.

This contrasts with policy-based approaches, which rely on runtime enforcement of delegation rules, and with mutable delegation architectures, which permit dynamic re-parenting without structural constraints. BX3's immutable parent chain is not a restriction on capability---the Bounds Engine still permits arbitrary task decomposition and subtask execution within the inherited scope---it is a guarantee about what happens when capability fails or accountability must be traced.

The framework's Layer model also provides a natural integration point for the protocols described in the preceding subsections. HDP's cryptographic tokens can encode BX3 layer boundaries as capability scopes. IBCTs can bind to BX3 Purpose Layer identifiers rather than to agent identities, enabling provenance chains that survive agent replacement. DCC's seven properties align with BX3's layer definitions: authority narrowing maps to Purpose Layer escalation, forensic reconstructibility maps to the immutable parent pointer, scope-action conformance maps to Bounds Engine constraint checking, and output schema conformance maps to Fact Layer deterministic validation.

The convergence of BX3's Layer model with these independently developed protocols suggests that the three-layer functional decomposition is not an arbitrary design choice but a natural structural pattern that emerges from the requirements of accountable, auditable, and non-deterministic-resilient multi-agent systems---and that recursive spawning, implemented on this substrate, provides the strongest available guarantee of upstream accountability in hierarchical agent architectures.